% !TeX root = ../../Book3.tex
\section{横割与素端}
\label{b3:sec:0701}

内部共形坐标不自动给出边界坐标。裂缝两侧在平面中是同一组点，却对应不同的趋近方式。横割与素端用区域内部的分离来记录这种方式。本节先在有界单连通区域 \(U\) 中定义；无界情形在球面中取边界、闭包与直径，允许无穷远作为端点；这给出相应的不依赖有限坐标的版本。

\subsection{横割}
\label{b3:subsec:070101}

\begin{definition}\label{b3:def:07-crosscut}
区域 \(U\) 的\term[hengge]{横割}[crosscut]是一个开 Jordan 弧 \(C\subset U\)，它的闭弧有两个不同的端点在 \(\partial U\)，其余点都在 \(U\)。
\end{definition}
对于单连通平面区域，横割把 \(U\) 分成恰好两个连通分支。这是 Jordan 分离定理在横割上的形式：沿区域内路径把弧的两侧闭合，得到一个 Jordan 圈；两侧不能绕过边界端点相通。圆盘中的弦是基本例子，弧可以弯曲而不必是直线。若弧的两个端点在内部，它通常不分离区域；若闭弧自交，也不能套用横割分离结论。

\ProseParagraphBreak
固定观察点 \(a\in U\)。当 \(a\notin C\) 时，记不含 \(a\) 的横割分支为 \(D(C)\)。它代表从 \(a\) 看被横割挡住的一侧。边界接近的定义将使用这些内部开集，而不是先假定每个边界点都有一条可达路径。

\ProseParagraphBreak
横割的两个端点属于边界，内部属于区域；这三项条件都参与分离结论。单位盘中的弦是最简单的模型：直弦把圆盘分成两个半盘型区域，圆周上两个开弧分别贴近两个分支。若把弦改成只以一个端点落在边界的线段，它的内部终点可以绕过去，不再把圆盘分开。若把两端都落在区域内部，也没有相同的两侧分离结论。一般单连通域的横割分离是平面 Jordan 拓扑的应用，并不是单凭曲线连续就能推出。

\ProseParagraphBreak
选定观察点的作用是给横割的两侧作一致选择。若观察点落在横割上，先把观察点换到区域内另一个点；链趋近边界后，两点最终处在同一个远侧分支，因此不会改变链的端部等价类。这里的“最终”很关键：前面有限几条横割的选边可以不同，但素端由链的尾部决定。把有限项删除或重新编号，不应制造新的素端。

\subsection{链}
\label{b3:subsec:070102}

\begin{definition}\label{b3:def:07-null-chain}
一列横割 \(C_n\) 称为一条零链，若闭横割两两不交，\(\operatorname{diam}C_n\rightarrow0\)，且在舍去有限项后，\(D_n=D(C_n)\) 满足 \(D_{n+1}\subset D_n\)、\(C_{n+1}\subset D_n\)，同时 \(\bigcap_nD_n=\varnothing\)。
\end{definition}
小直径约束的是横割而非整个 \(D_n\)。复杂边界中一个很短的横割后面仍可能藏着一个很长的狭窄通道。链的印象定义为
\[
 I=\bigcap_n\overline{D_n}.
\]
在有界情形，它是非空紧连通集且包含于边界。非空来自嵌套紧集，连通性来自连通闭包的嵌套交；若交内有内点，则该点的小圆盘与充分小、远离此点的横割不交，连通性迫使它一直处在同一侧，最终与内部交为空矛盾。不能仅凭横割直径趋于0就把印象写成单点。

\ProseParagraphBreak
对于圆盘，固定边界点附近的短弦给出一条链，分支为越来越小的圆帽，印象恰为该点。对于裂盘，同一个裂缝内点有从上方与下方趋近的两条链；它们的印象可能相同，却不代表同一个边界通道。

\subsection{素端}
\label{b3:subsec:070103}

若对每个 \(n\)，另一条链的充分晚分支都包含在 \(D_n\) 中，称后一链从属前一链。两链互相从属时称为等价。传递性由“先给第一链的一个分支，再依次选择后两链的充分晚分支”得到，所以等价是等价关系。

\begin{definition}\label{b3:def:07-prime-end}
零链的等价类称为 \(U\) 的\term[suduan]{素端}[prime end]。其印象是任一代表链的闭分支交。
\end{definition}
互相从属保证两个代表给出的印象相同。改变观察点也不影响充分晚的从属关系：连接两个观察点的固定内部紧路径不会与充分晚的横割相交，两个点最终落在同一侧。定义依靠尾部，故舍掉有限条横割不改变素端。

\ProseParagraphBreak
“素端”与“边界点”属于不同集合。印象是一个集合，素端是趋近通道的等价类；既可能一个素端有非单点印象，也可能不同素端有相同单点印象。对有规则边界的区域，二者恰巧一致；这需要边界对应定理，不能写入一般定义。

\subsection{素端边界}
\label{b3:subsec:070104}

令 \(\widehat U^{\mathrm P}=U\cup\partial_{\mathrm P}U\)。内部沿用通常拓扑；由一个链分支 \(D_n\) 加上所有从属于该分支的素端，定义该素端的一组邻域。互相从属使这种拓扑与代表链无关。这一空间称为素端紧化，其边界是 \(\partial_{\mathrm P}U\)。

\begin{theorem}[深层素端对应定理]\label{b3:thm:07-prime-end-correspondence}
设 \(U\subsetneq\mathbb C\) 单连通，\(f:\mathbb D\rightarrow U\) 双全纯。则 \(f\) 唯一延为 \(\overline{\mathbb D}\) 到 \(\widehat U^{\mathrm P}\) 的同胚。特别地，素端边界具有一个圆周的拓扑与循环次序。
\end{theorem}
\begin{proofsketch}
这个定理包含一般平面边界的分离与链逼近，不在这里声称完成全部拓扑证明。路线是用横割分支的嵌套关系建立素端拓扑；再用共形映射的面积估计，选取趋向圆周的短弧，使其像形成零链；反向以圆盘内的横割分离证明同一素端只能对应一个圆周点；最后核对分支邻域对应圆帽邻域。困难在于普通边界可能并非局部连通，不能假定横割像的闭弧已有边界端点。完整处理可读 Lars V. Ahlfors 的《Conformal Invariants》\cite{Ahlfors2010Conformal}的 \S4-6（57--60页），其中也说明横割链与圆盘边界的对应；Christian Pommerenke 的《Boundary Behaviour of Conformal Maps》\cite{Pommerenke1992Boundary}第2章则为进一步阅读。这里保留该定理为明确的深层输入。
\end{proofsketch}
这个输入给边界一个合适的拓扑对象，并没有保证 \(f\) 在每个圆周点有唯一普通边界极限。由素端印象到平面边界的投影可能有重合或非单值集合像。下节说明何时能把这项结果化为通常的连续延拓。

\ProseParagraphBreak
裂缝圆盘提供必须保留两侧的例子。把单位盘沿正实半径切开，在裂缝内部的一个边界点 \(x\in(0,1)\)，可以从上侧和下侧分别取越来越小的半圆横割。两条链的印象都是 \(\{x\}\)，但选出的近侧区域不最终互相包含：一条始终位于裂缝上侧，另一条始终位于下侧。因此它们对应两个不同素端。素端记录趋近方向，印象只记录最后落在哪些平面边界点，二者不能一一混同。裂缝端点附近的链又可能把两侧接在同一个角端，须看局部横割的实际分离。

\ProseParagraphBreak
在圆盘中，边界点 \(\zeta\) 的标准链可取 \(\mathbb D\cap\{|z-\zeta|=\varepsilon_n\}\)，其中 \(\varepsilon_n\downarrow0\)。选靠近 \(\zeta\) 的分支，其闭包交为 \(\{\zeta\}\)。链邻域与通常的圆盘闭包邻域相容。映射到一般单连通域后，平面直径未必仍按同样速度趋于0，素端对应定理正是把这种内部对应和边界拓扑协调起来；它的证明不能仅用连续映射将小圆变成小曲线的一句话代替，因为此时边界连续性尚未建立。
